\section{Introduction}\label{sec:intro}

This paper reconstructs the argument for the following theorem.

\begin{theorem}\label{thm:intro}
For every primitive Dirichlet character $\chi$, the meromorphic function $L(s,\chi)$ has no zeros in the half-plane $\operatorname{Re}s>\frac12$. In particular, the Riemann zeta function has no zeros in $\operatorname{Re}s>\frac12$.
\end{theorem}

The proof is due to Levent Alpöge and Claude. On November 20, 2026, Alpöge released a four-page note~\cite{note} together with a proof of Theorem~\ref{thm:intro} formalized in Lean~\cite{lean4} on top of Mathlib~\cite{mathlib}, and on November 21, 2026, the Lean and Mathlib communities, with the endorsement of the Lean FRO and the Mathlib Initiative, certified that the formal statement and the formal proof are correct~\cite{leanstatement}. The note describes the argument as a descendant of Stepanov's method~\cite{stepanov} for curves over finite fields. It credits Alpöge with the plan and with the recognition that an auxiliary-function argument of Siegel, which appeared in mathematical work by researchers at OpenAI, is Stepanov's method with multiplicity one. Alpöge describes the turning point as the appearance of multiplicity two in an automorphic determinant on $\mathrm O(2,n)$, which suggested that further theta lifts could increase the multiplicity. The note credits Claude with the proposal that a chain of theta lifts of length about $\sqrt x$ should play the role of multiplicity, and with the arguments on automorphic forms on groups of higher rank. It thanks David Dohan, Ben Michel, Ralph Furman, Eric Price and Cora Vinther. This paper writes the proof as conventional mathematics that a reader can check without the formal development. The reconstruction is not complete. The paper proves Theorem~\ref{thm:intro} from the published theorems it cites and from a finite list of internal statements, each stated precisely where it is used and summarized in Subsection~\ref{sec:frontier}, whose conventional proofs it does not yet give. The formal development proves each of these statements. A few of the cited published theorems have not yet been compared line by line with their sources, and Subsection~\ref{sec:frontier} lists them as well.

\subsection{The architecture of the proof}

For $x\ge1$ let
\[
G(x)=\sum_{n\le x}\Lambda(n)\Big(\frac xn-1\Big)-x\log x+(1+\gamma)x ,
\]
and let $G_\chi$ be the analogous function for a Dirichlet character $\chi$, defined in \eqref{Gchi}. The Mellin transform of $G$ has a simple pole at each zero of $\zeta$ in the critical strip, with residue $-1/(\rho(\rho-1))$ times the multiplicity. The proof has two parts.

\emph{Part I} reduces Theorem~\ref{thm:intro} to one arithmetic statement. Section~\ref{sec:criterion} proves a quantitative form of Landau's theorem on Dirichlet integrals of one sign. If a positive combination $G+\sum_iw_iG_{\chi_i}$ is at most $Cx^\theta$ for every $x$ in windows $Y\le x\le Y^{3/2}$, for arbitrarily large $Y$, then $\zeta$ and the $L(s,\chi_i)$ do not vanish for $\operatorname{Re}s>\theta$. So it suffices to prove such a bound with $\theta=\frac12+O(\varepsilon)$, on windows, for the powers of a fixed primitive character $\psi_0$. The bound comes from an abelian CM field $E$ of degree $4n$ with $n$ about $x^{2+8\varepsilon}$, cut out by $\psi_0$ and two auxiliary primes $q$ and $d$. Section~\ref{sec:accounting} proves two identities. The first, Proposition~\ref{acct}, shows that seven inequalities between slopes of hermitian lattices, heights of period maps and local terms bound a sum $S$ over the primes inert in $\mathbb Q(\sqrt{-d})$ by an explicit main term. The second, Proposition~\ref{excess}, identifies $S$ minus that main term with a combination of the functions $G_\psi$ over the characters $\psi$ of $\mathrm{Gal}(E/\mathbb Q)$. Section~\ref{sec:point} turns a bound for this combination, summed over many weights and many nodes, into a bound for a positive combination of the $G_\psi$ at one point $x$. An average over translations removes the characters outside a small family, and an equidistributed choice of weights gives every remaining character a positive weight. The step leaves the values $\lvert G_\psi(x)\rvert$ for $3o$ auxiliary characters, where $o$ is the order of $\psi_0$. Section~\ref{sec:auxiliary} chooses $q$ and $d$ by a mean square over families of prime moduli so that these values are at most $x^{1/2+\varepsilon}$ on a whole window, without assuming any zero-free region beyond the classical one, and constructs the field $E$ and the remaining data. Section~\ref{sec:reduction} states the arithmetic theorem, Theorem~A, and deduces Theorem~\ref{thm:intro} from it.

\emph{Part II} concerns Theorem~A. Bombieri's form of Stepanov's argument~\cite{bombieri} bounds the number of rational points of a curve by comparing an auxiliary function of controlled pole order with its order of vanishing $M$ at the rational points, and the two errors balance at $M=\sqrt q$. Over a number field the count of zeros against poles becomes a slope inequality of Bost's type~\cite{bost} for a family of maps between hermitian lattices, and the order of vanishing becomes the number $m$ of steps in a chain of theta lifts between unitary groups of growing rank. To each step the argument attaches a hermitian lattice over the ring of integers of $E$ and a period map into the dual of a lattice of automorphic forms. At a prime $p$ inert in $\mathbb Q(\sqrt{-d})$, an elliptic curve with complex multiplication by $\mathbb Q(\sqrt{-d})$ has a formal group of height two, and the theta kernel of one lift is divisible by a power of $p$ whose exponent at weight $\alpha$ has the main term $\sum_{i\ \mathrm{odd}}\alpha/p^i=\alpha p/(p^2-1)$. This main term matches the contribution of $p$ to the sum $S$. Hecke operators at the inert primes cost a loss that a nondegeneracy theorem for integral representations of $\mathrm{GL}_k$ controls. The archimedean size of the period map is $m$ times a quantity close to $\sum_t\log\Gamma(a_t)$, measured against the period of the CM curve, and arithmetic intersection theory on unitary Shimura varieties bounds the slope of the lattice of automorphic forms. Bessel Poincar\'e series on a unitary group realize the period maps at one complex place, and their near orthogonality gives the slope inequality for the family. At the places where the local type of the forms is twisted, lengths of Tate--Shafarevich groups of the CM field bound the congruence defects that enter the heights, through Galois representations attached to automorphic representations of unitary groups and $p$-adic Hodge theory. The two error terms balance at $m=\sqrt x$, as in Stepanov's argument. Section~\ref{sec:organization} describes the sections of Part~II.

\subsection{Contemporaneous accounts}

Several groups began translating, mapping, explaining or distilling the formal development within days of its release. Zhang~\cite{zhang} wrote a preliminary mathematical account of the proof and of its principal ingredients, with an emphasis on their sources and history. Rubinstein~\cite{rubinsteinpage,rubinsteinzeta} maintains a guide to the proof together with a detailed exposition of the argument specialized to the Riemann zeta function, and his verification run compiled the formal development independently. Naskręcki~\cite{naskrecki} built a map of the formal proof, its structure and dependencies, and a record of its verification status, together with an atlas of the route. Epstein~\cite{epstein} wrote a human-readable blueprint that identifies the top-level statements, separates the principal mathematical crux from library material, and draws the dependency structure. Brox~\cite{brox} leads a program to distill and simplify the formal proof into shorter human-readable proofs of the Riemann hypothesis and of the generalized Riemann hypothesis for Dirichlet characters, which is in progress. Slijepčević~\cite{slijepcevic} wrote an outline of the proof in standard language, including the positivity argument, the theta-lift and lattice estimates, the arithmetic losses and the Galois-representation input. The present paper writes the argument as one continuous conventional proof with its open steps stated, so that a reader can judge its proofs and its cited external inputs without appeal to the formal development as mathematical authority.

\subsection{How this account was obtained}

The formal development has about $1.07$ million declarations, of which about $767{,}000$ lie on the dependency route of the final theorem, and a readable proof cannot arise by translating them one at a time. We began at the formal theorem \texttt{riemannHypothesis}, traced its dependencies backward, and identified the mathematical statements on which it depends. Whenever a region of the development implemented a recognizable argument, we stated and proved that argument at its natural mathematical level, and whenever the route reached an established theorem of the literature, we cited that theorem and did not use the formal development as its justification. This process does not repeat the formal verification of~\cite{leanstatement}, and it did not involve reading every declaration. The author used Claude (Fable 5.2 and Opus 5.5), GPT (including GPT 5.6 Pro) and Kimi K3.1 Pro in the reconstruction. The formal certification shows that the Lean term proves the formal statement. The reconstruction writes the content of the internal dependencies as mathematics, and a separate comparison of each cited theorem with its published statement checks hypotheses, normalizations and quantifiers.

Two findings of the reconstruction concern the strength of the internal statements. The bound for Tate--Shafarevich groups on windows of scales, Hypothesis~\ref{w4:thm:vwin}, implies a bound in the direction of Vandiver's conjecture that is stronger than the published bounds of Soul\'e~\cite{soule99} and of de~la~Torre~Dur\'an~\cite{delatorre} (Proposition~\ref{w4:prop:need}). One branch of the formal proof uses an automorphy statement for residually reducible Galois representations of unitary type in three variables, (G3) of Subsection~\ref{tw:sec:vandiver}, which lies outside the hypotheses of the published automorphy lifting theorems for residually reducible representations. Second, the bound for the covolumes of the lattices of automorphic forms needs a constant independent of the level of a deep cover of a unitary Shimura variety, and the argument obtains it from the arithmetic Hilbert--Samuel theorem on each model of a tower of models and an exact identity at the good primes, without a uniform version of the Hilbert--Samuel theorem (Subsection~\ref{w3:sec:unif}).

\subsection{Organization}\label{sec:organization}

Sections~\ref{sec:criterion} to~\ref{sec:reduction} form Part~I. Part~II begins in Section~\ref{sec:chain} with the theta kernels between hermitian lattices over the CM field and their covariance. Section~\ref{sec:inert} studies the theta function of a CM elliptic curve at an inert prime and proves the form of the kernels in the parameters of a formal group of height two. Section~\ref{sec:gain} classifies the semilinear pairings on that formal group and proves the valuation estimate for the kernels. Section~\ref{sec:chaindiv} turns this estimate into divisibility of the composite of a chain of lifts. Section~\ref{sec:lifts} constructs the lifts and shows that they form a chain datum. Section~\ref{r2:sec:loss} proves the nondegeneracy of Hecke operators that controls the loss, and hypotheses \eqref{A1} and \eqref{A2}. Section~\ref{r3:sec} proves the slope inequality \eqref{A3} for the family of period maps, and Section~\ref{sec:bessel} proves its automorphic inputs for Bessel Poincar\'e series. Section~\ref{r4:sec:R4} computes the period of the CM curve and proves \eqref{A4}, \eqref{A6} and \eqref{A7}. Section~\ref{r5:sec:localterms} proves \eqref{A5}. Section~\ref{sec:twisted} treats the twisted places, the Tate--Shafarevich bound and the congruence bound. Section~\ref{sec:proofA} constructs the data, the weights and the nodes, and proves Theorem~A from the statements of the preceding sections. Subsection~\ref{sec:frontier} collects the hypotheses used in the final deduction.
